\section{Realizable Karp Factor}

The bounded-cycle audit separates abstract cycles in the finite tail graph from
cycle words that lift to positive integer orbits.  In the joint Karp/slope
sweep at \((q,R_{\max})=(5,4),(6,5),(7,6)\), with the automaton period widened
across the levels, the only slope-compatible survivor at each tested level is
the elementary valuation word \([2]\), with factor \(3/4\)
\(\artifact{karp_slope_joint_sweep.json}{karp\_slope\_joint\_sweep.json}\).

The realizability extension audits all simple cycles up to
\(\texttt{max\_cycle\_edges}=12\) at the same three levels.  It classifies
\(226333\) simple cycles: \(226330\) are
\(\texttt{noninteger\_2adic\_only}\), and the only \(3\)
\(\texttt{positive\_integer\_cycle}\) entries are the elementary \([2]\) cycle,
one per level
\(\artifact{tail_cycle_realizability.json}{tail\_cycle\_realizability.json}\),
\(\artifact{tail_cycle_realizability_extended.json}{tail\_cycle\_realizability\_extended.json}\).
Thus the saved finite diagnostic has
\[
  \text{realizable Karp factor}
  =
  \text{slope-filtered Karp factor}
  =
  \frac34
\]
on these levels.

This is a finite bounded-cycle fact, not a global statement about arbitrary
infinite products.  The high-growth abstract cycles in the LTE-closed graph are
not thereby eliminated from the infinite problem; the audit only records that,
within the tested bounded simple-cycle regime, their cycle words do not lift to
positive integer cycles.

The value \(3/4\) should also be distinguished from the modular-transition
number appearing in Chang's human-LLM Collatz manuscript
\cite{chang2026collatz}.  In this framework it is a cycle-mean factor for an
elementary valuation word, not a modular transition probability.
